\section{Conclusion}
\label{sec:conclusion}

We presented a method that reconstructs road collisions from independent
onboard recordings by connecting evidence across recorders through clock
recovery, physical vehicle association and retrospective motion
estimation. This includes observations of vehicles without their own
logs. Through semantic conversion, the resulting chronology relates
manoeuvres and conflicts to recorded control actions and contacts while
keeping source observations and the evidence behind identity decisions
traceable.

In the reference condition, with impact reports from both participants,
the method correctly identified all 26 collision participant sets and
represented all 45 observable
non-recording vehicles across 19 simulated scenarios. Across the
1816-job evaluation, tested positioning and nominal radar errors
preserved collision identification and high local conflict precision.
When every recorder underwent extreme radar corruption, 95.4\% of
collision participant sets remained correct while local conflict
precision fell to 31.3\%, revealing the limit of semantic interpretation
under dense clutter. This chronology supports investigators in
examining competing accounts of a collision.

Future work will extend validation to real onboard recordings and
examine how motion and identity uncertainty can inform confidence
estimates for semantic events under dense radar clutter. We will also
investigate large language models for abductive reasoning over kinematic
and semantic traces, generating and comparing candidate collision
explanations against supporting and contradicting recorded evidence.
